% from ega4-8-fr.tex
\oldpage[IV-3]{53}
Supposons d'abord que $X$ soit aussi affine~; l'assertion à prouver est alors
équivalente au

\begin{corollary}[8.14.2.2]
\label{IV.8.14.2.2-fr}
Soient $A$ un anneau, $B$ une $A$-algèbre. Afin que, pour tout système
inductif filtrant $(C_\lambda)$ de $A$-algèbres, l'application canonique
\begin{equation}
\label{IV.8.14.2.3-fr}
\varinjlim_\lambda\operatorname{Hom}_{A\text{-}\mathrm{alg.}}(B,C_\lambda)
\longrightarrow
\operatorname{Hom}_{A\text{-}\mathrm{alg.}}
(B,\varinjlim_\lambda C_\lambda) \tag{8.14.2.3}
\end{equation}
soit bijective, il faut et il suffit que $B$ soit une $A$-algèbre de
présentation finie.
\end{corollary}

Il reste seulement à montrer que la condition est nécessaire. Prenons d'abord
pour $(C_\lambda)$ le système inductif filtrant des sous-$A$-algèbres de type
fini de $B$, de sorte que $\varinjlim_\lambda C_\lambda=B$. Le fait que (8.14.2.3)
soit bijective entraîne en particulier que l'application identique
$1_B$ se factorise en $B\to C_\lambda\to B$ pour un $\lambda$ convenable,
ce qui entraîne $C_\lambda=B$, donc $B$ est une $A$-algèbre \emph{de type
fini}. Posons alors $B=C/\mathfrak J$, où
$C=A[T_1,\ldots,T_n]$ et $\mathfrak J$ est un idéal de $C$. Alors
$\mathfrak J$ est limite inductive filtrante des idéaux de type fini
$\mathfrak J_\lambda\subset\mathfrak J$ de $C$~; posant
$C_\lambda=C/\mathfrak J_\lambda$, et utilisant l'exactitude du foncteur
$\varinjlim$, on voit que $B$ est encore isomorphe à la limite inductive du
système inductif filtrant $(C_\lambda)$. Il existe donc un $\lambda$ et un
$A$-homomorphisme $u:B\to C_\lambda$ tels que le composé
\[
B\xrightarrow{u}C_\lambda\xrightarrow{p_\lambda}B
\]
(où $p_\lambda$ est l'homomorphisme canonique) soit l'identité. Soit
$q_\lambda:C\to C_\lambda$ l'homomorphisme canonique, et posons
$t_i=p_\lambda(q_\lambda(T_i))$~; on a donc
$p_\lambda(u(t_i))=p_\lambda(q_\lambda(T_i))$, autrement dit
\[
u(t_i)-q_\lambda(T_i)\in\mathfrak J/\mathfrak J_\lambda.
\]
Il existe donc un $\mu\geqslant\lambda$ tel que les $n$ éléments
$u(t_i)-q_\lambda(T_i)$ appartiennent à
$\mathfrak J_\mu/\mathfrak J_\lambda$ $(1\leqslant i\leqslant n)$~; si
$p_{\mu\lambda}:C_\lambda\to C_\mu$ est l'homomorphisme canonique, on a par
suite
\[
p_{\mu\lambda}(u(t_i))
=p_{\mu\lambda}(q_\lambda(T_i))=q_\mu(T_i).
\]
Remplaçant $\lambda$ par $\mu$ et $u$ par $p_{\mu\lambda}\circ u$, on peut
donc supposer que $u(t_i)=q_\lambda(T_i)$ pour tout $i$, et si
$r=p_\lambda\circ q_\lambda$ est l'homomorphisme canonique
$C\to C/\mathfrak J=B$, on peut donc écrire
$u(r(T_i))=q_\lambda(T_i)$ pour tout $i$, d'où
$q_\lambda=u\circ r$. Mais cela entraîne nécessairement que
$\mathfrak J=\mathfrak J_\lambda$, car si $z\in\mathfrak J$, on a
$r(z)=0$~; donc on a $B=C_\lambda$.

Passons maintenant au cas où $S$ est affine et $X$ quelconque~; tout revient à
prouver qu'un ouvert affine $V$ de $X$ est de présentation finie sur $S$, et
en vertu de ce qui vient d'être démontré, il suffit de prouver que pour tout
système projectif filtrant $(Z_\lambda)$ de $S$-préschémas affines,
l'application
\begin{equation}
\label{IV.8.14.2.4-fr}
\varinjlim_\lambda\operatorname{Hom}_S(Z_\lambda,V)
\longrightarrow
\operatorname{Hom}_S(\varprojlim_\lambda Z_\lambda,V) \tag{8.14.2.4}
\end{equation}
est bijective. Il est immédiat que cette application est injective, car si
$(v_\lambda)$, $(v'_\lambda)$ sont deux systèmes inductifs de
$S$-homomorphismes $v_\lambda:Z_\lambda\to V$,
$v'_\lambda:Z_\lambda\to V$ tels que les morphismes correspondants
\[
Z\xrightarrow{u_\lambda}Z_\lambda\xrightarrow{v_\lambda}V,
\qquad
Z\xrightarrow{u_\lambda}Z_\lambda\xrightarrow{v'_\lambda}V
\]
soient égaux ($u_\lambda$ étant le morphisme canonique), alors les morphismes
\[
Z\xrightarrow{u_\lambda}Z_\lambda\xrightarrow{v_\lambda}V
\xrightarrow{j}X,
\qquad
Z\xrightarrow{u_\lambda}Z_\lambda\xrightarrow{v'_\lambda}V
\xrightarrow{j}X
\]
(où $j$ est l'injection canonique) sont égaux, ce qui entraîne
$j\circ v_\lambda=j\circ v'_\lambda$ par hypothèse pour un $\lambda$
convenable, donc $v_\lambda=v'_\lambda$.

